\section{Related Work}
\label{sec:related}

\subsection{UAV communication and FANET systems}

UAV-enabled wireless communication has been widely studied because aerial platforms can provide on-demand coverage, relay support, flexible topology control, and rapid deployment in infrastructure-limited scenarios \cite{zeng2016uav,mozaffari2019tutorial}. Channel geometry, altitude, energy constraints, mobility, and interference are central design factors. The altitude-dependent coverage model in \cite{alhourani2014optimal}, for example, illustrates that UAV placement has a direct effect on link quality. FANET research further considers multiple UAVs with fast topology changes and mission-oriented networking constraints \cite{lakew2020fanet}. Within the target journal, energy-efficient FANET routing \cite{oubbati2019ecad}, mission-oriented UAV swarm fault tolerance \cite{wang2021fault}, and UAV simulation tools \cite{jiang2024simulation} show that communication-system-level modeling and reproducible simulation remain important research directions.
Recent IJCS work has also considered UAV-assisted anti-jamming and proactive handoff with deep-learning support \cite{unnisa2023antijamming}, which is close to the present paper's emphasis on adversarial wireless environments.

The present work is related to this literature but has a different focus. We do not study multi-hop FANET routing or swarm task allocation. Instead, we study a secure communication control problem in which two cooperative UAVs must serve users and jam an eavesdropper while maintaining a synchronized digital twin of the adversarial state.

\subsection{Physical-layer security for UAV communications}

Physical-layer security originates from information-theoretic secrecy models such as the wiretap channel \cite{wyner1975wiretap} and broadcast channels with confidential messages \cite{csiszar1978broadcast}. In UAV systems, secrecy depends heavily on geometry because UAV mobility changes the legitimate and eavesdropping channel gains. Joint trajectory and power optimization has therefore become a common approach. Zhang et al. \cite{zhang2019securing} optimized UAV trajectory and transmit power for secure UAV communication. Cooperative jamming has also been introduced to improve secrecy; Zhong et al. \cite{zhong2019cooperative} studied a two-UAV setting with a transmitter UAV and a jammer UAV, while Zhou et al. \cite{zhou2018unknown} considered a friendly UAV jammer when the eavesdropper location is uncertain. Sun et al. \cite{sun2019aerial} further examined aerial cooperative jamming with joint trajectory and power control.

Recent work has moved toward more realistic secure UAV settings with sensing, edge computing, channel uncertainty, and learning-based optimization. Wei et al. \cite{wei2024isncsecureuav} considered an integrated sensing, navigation, and communication framework for a secure UAV network with a mobile aerial eavesdropper. Mao et al. \cite{mao2023secureuavmec} and Lu et al. \cite{lu2024securemaritimemec} studied secure UAV-assisted MEC systems with coupled trajectory, resource, and security constraints. Ding et al. \cite{ding2024ddqnuavmec} used a DDQN-based method for trajectory and resource optimization in UAV-aided MEC secure communications, while Liu et al. \cite{liu2024secureisacuav} and Chen et al. \cite{chen2025covertuavmec} addressed recent ISAC-UAV and covert UAV-MEC security problems.

These works motivate our use of trajectory, transmit power, and jamming power as control variables, and they provide the most relevant recent context for the baselines in this paper. The key distinction is the information state. Rather than assuming a fixed eavesdropper state, a one-shot uncertainty set, or an MEC/offloading objective, we model a digital twin whose quality evolves according to synchronization decisions. Consequently, secrecy performance depends not only on the physical trajectory and power allocation, but also on whether the controller has recently synchronized the twin. Directly importing the DDQN or SCA implementations from these papers would change the state, action, constraints, and reward definitions; we therefore compare against matched in-simulator rollout, rule-based, and SCA-style diagnostic baselines, and we state this scope explicitly in the experiments.

\subsection{Digital twins and information freshness}

Digital twins have been proposed as a means to support intelligent wireless systems, network automation, and data-driven control. Surveys on 6G and wireless digital twins discuss their role in modeling physical network behavior, supporting closed-loop control, and reducing the risk of deploying untested policies \cite{kuruvatti2022digital}. For UAV networks, digital-twin-enabled domain adaptation has been identified as a promising direction for zero-touch control, simulation-to-real transfer, and robust policy evaluation \cite{mcmanus2023digital}.
Recent digital-twin communication work is closer to the present setting. Li et al. \cite{li2023adaptivedt} studied adaptive digital twins for UAV-assisted integrated sensing, communication, and computation networks. Yang et al. \cite{yang2024syncdelay} optimized synchronization delay for wireless digital twins, and a related joint communication-computation framework for digital twins over wireless networks was developed in \cite{yang2024jointdt}. These papers emphasize that synchronization and communication/computation resources must be optimized jointly, which aligns with the motivation of this paper.

The synchronization aspect of our problem is also related to age of information (AoI), which measures the freshness of status updates. Early AoI work showed that update frequency should be optimized rather than simply maximized, because network resources and queueing effects create tradeoffs \cite{kaul2012aoi}. Vehicular-network AoI studies reached similar conclusions for mobile systems \cite{kaul2011vehicular}. Our setting shares the same core idea that stale information is harmful, but the performance target is different: we optimize secrecy rate and secrecy outage under a limited synchronization budget, not only information freshness.

\subsection{Optimization, reinforcement learning, and empirical risk estimation}

Many UAV secrecy problems are nonconvex because channel gains, trajectory variables, and power decisions are coupled. Successive convex approximation and alternating optimization are therefore widely used in recent secure UAV optimization studies \cite{mao2023secureuavmec,wei2024isncsecureuav,lu2024securemaritimemec,liu2024secureisacuav}, with convex optimization providing the standard mathematical foundation \cite{boyd2004convex}. In this paper, the SCA-based baselines are matched local-optimization diagnostics under either deployable twin information or oracle eavesdropper information; they are not presented as a reproduction of a particular external paper, and Boyd's text is cited only as background for the convex-optimization machinery.

Reinforcement learning is another natural approach for sequential control. PPO \cite{schulman2017ppo} is a widely used policy-gradient method and is included as a lightweight learning-based diagnostic in our experiments. We do not claim that PPO is intrinsically unsuitable; our result only evaluates a limited-training run under the same simulator.

Finally, our secrecy-loss risk estimator uses a simple residual-calibration idea: fit a loss model, add a calibration residual quantile, and validate the resulting upper bound on held-out simulator traces. We do not present it as a contribution to formal predictive inference, nor do we claim a finite-sample coverage theorem under arbitrary deployment conditions. It is a holdout-fitted empirical upper bound that is calibrated and validated on simulated trajectory distributions. This conservative interpretation is important for avoiding overstatement of the estimator's scope.
